% 362_BAW_Ising_FFGFT_De_ch.tex
% Johann Pascher, 12. Sept. 2026
% Die BAW-Ising-Maschine (Vadde et al., Commun. Phys. 9:290, 2026)
% als Resonanzrechner im Sinne der FFGFT

\providecommand{\BM}{\textbf{[B]}}
\providecommand{\KM}{\textbf{[K]}}
\providecommand{\SM}{\textbf{[S]}}
\providecommand{\SetzM}{\textbf{[SETZUNG]}}
\providecommand{\QM}{\textbf{[Q]}}

\chapter*{Die BAW-Ising-Maschine als Resonanzrechner}
\addcontentsline{toc}{chapter}{Die BAW-Ising-Maschine als Resonanzrechner}

\begin{abstract}
Vadde, Ovcharov, González, Khymyn, Litvinenko und Åkerman (Commun.\
Phys.~9:290, 8.~September 2026, Open Access) stellen eine
zeitmultiplexte Ising-Maschine mit 2048 Spins auf zwei Quarz-Volumen\-schall\-wellen-%
Verzögerungsleitungen vor. Dieses Dokument liest die Arbeit als
Fallstudie der FFGFT-Einordnung von Resonanzrechnern (Dok.~173, 328, 343):
Spins als Phasenzustände eines Ringoszillators, Barkhausen-Bedingung als
Wicklungsquantisierung, parametrische $\mathbb{Z}_2$-Selektion, Kopplungsregime
mit Chaosschwelle, und -- quantitativ am deutlichsten -- ein Auflösungsboden:
Näherungslösungen zu 99,9\,\% werden zuverlässig gefunden, exakte Lösungen
($E=0$) mit wachsender Problemgröße nicht mehr. Alle Zitate wörtlich \QM.
Status der Einordnungen: \SM.
\end{abstract}

% ============================================================
\section*{Statusmarkierungen}
% ============================================================
\begin{center}
\begin{tabular}{@{}lp{10cm}@{}}
\textbf{[B]} & Algebraisch bewiesen. \\[3pt]
\textbf{[K]} & Numerisch bestätigt. \\[3pt]
\textbf{[S]} & Strukturell plausibel, physikalische Identifikation offen. \\[3pt]
\textbf{[Q]} & Aus der Quelle (Vadde et al.) übernommen, wörtlich zitiert. \\[3pt]
\textbf{[SETZUNG]} & Axiomatische Setzung. \\
\end{tabular}
\end{center}

\noindent Alle wörtlichen Zitate stammen aus Vadde et al.\ (2026), Open Access
unter CC-BY~4.0; Abschnittsangaben beziehen sich auf die dortige Gliederung.

% ============================================================
\section{Die Quelle: Was Vadde et al.\ bauen}
% ============================================================

\subsection{Architektur}

\begin{quote}\QM
\enquote{We use two serially connected 20.5 MHz, 707~$\mu$s bulk acoustic
wave delay lines supporting 2048 spins. Our design provides all-to-all
connectivity with 15-bit coupling resolution and finds approximate MAX-CUT
solutions in 341~ms.} (Abstract)
\end{quote}
Die Spins sind keine Teilchen, sondern Phasen zirkulierender RF-Pulse:
\begin{quote}\QM
\enquote{The spin states $+1$ and $-1$ are encoded in the relative phase of
the pulses, with bistable phases achieved by phase-binarizing the
propagating RF pulses using a phase-sensitive amplifier (PSA).}
(Abschnitt \emph{Design})
\end{quote}

\subsection{Der Ringoszillator und die Barkhausen-Bedingung}

\begin{quote}\QM
\enquote{The BAW delay line, together with linear amplifiers and the PSA,
forms a multi-physics ring oscillator. The key step of the design is to
establish a stable circulation of RF pulses within this loop. This requires
satisfying the Barkhausen stability criteria: the loop gain must be larger
than one, and the total phase shift around the loop must be an integer
multiple of $2\pi$.} (Abschnitt \emph{Design})
\end{quote}

\subsection{Parametrische Binarisierung}

\begin{quote}\QM
\enquote{The PSA further narrows the stability criteria to only those
signals with either phase 0 or $\pi$ relative to the pumping signal that is
twice the reference frequency.} (Abschnitt \emph{Design})
\end{quote}

\subsection{Mindestzyklen und Pulsdauer}

\begin{quote}\QM
\enquote{Due to the resonance mechanism underlying phase-sensitive
amplification, a minimum of 5--10 cycles per pulse is required, setting a
lower bound on the pulse duration to $\sim$250--500~ns. This constraint
dominates and effectively determines the pulse period.}
(Abschnitt \emph{Design})
\end{quote}

\subsection{Kopplungsstärke und Chaosgrenze}

\begin{quote}\QM
\enquote{The amplitude of the coupling signal can be tuned using the
potentiometer of the DAC, allowing an additional control to ensure the
injected signal remains within 5--30\,\% of the circulating RF signal and
does not overpower it, avoiding chaotization of the oscillatory circuit.}
(Abschnitt \emph{Design})
\end{quote}
\begin{quote}\QM
\enquote{The performance of BAWIM depends on the coupling strength, and
there exists an optimal coupling strength depending on the problem
instance.} (Abschnitt \emph{Performance})
\end{quote}

\subsection{Näherung gegen Exaktheit}

MAX-CUT:
\begin{quote}\QM
\enquote{Here, we observe that BAWIM attains up to 95\,\% of HbSB's best
performance in terms of Ising energy, and up to 99.99\,\% in terms of the
MAX-CUT score.} (Abschnitt \emph{Dependence on the edge density})
\end{quote}
Number Partitioning:
\begin{quote}\QM
\enquote{Here, the exact solution corresponds to a difference of 0, while
the 99.9\,\% approximate solution indicates that the difference is less than
0.1\,\% of the maximum possible value [\ldots]. In all cases, BAWIM
outperforms the HbSB algorithm, achieving a 100\,\% success rate in most
instances when considering the 99.9\,\% approximate solution.}
(Abschnitt \emph{Number partitioning})
\end{quote}
Die zugehörige Abbildung~5(b) zeigt, was der Text nicht ausspricht: Die
Kurve \enquote{BAWIM Exact} fällt von etwa 90\,\% Erfolg bei 32 Zahlen auf
nahe null bei 2048 Zahlen, während \enquote{BAWIM 99.9\,\%} über den
gesamten Bereich bei 100\,\% bleibt.

Sudoku:
\begin{quote}\QM
\enquote{Notably, even when the Ising energy approaches the ground state,
the Sudoku solution remains far from correct. [\ldots] Sudoku requires
exact constraint satisfaction; therefore, a near-ground-state Ising energy
may still correspond to an invalid solution.} (Abschnitt \emph{Sudoku})
\end{quote}

\subsection{Thermische Stabilität}

\begin{quote}\QM
\enquote{In contrast, our BAWIM exhibits a temperature coefficient of phase
accumulation of only 780~deg/$^\circ$C, which is $2.21\times10^4$ times
lower than that of the CIM.} (Abschnitt \emph{Comparison with CIMs})
\end{quote}

% ============================================================
\section{Einordnung 1: Ising-Maschine = Resonator (Dok.~173)}
% ============================================================

Dok.~173 §\emph{The Arc Closes} sagt über die kohärente Ising-Maschine:
\enquote{This resonator also has discrete states: it oscillates in either one
or the other phase -- spin-up or spin-down, zero or one. The excitation is
analogue -- continuous laser light, continuously adjustable coupling
strengths. The state is discrete.} Und in der Schlussfolgerung:
\enquote{An Ising machine is a resonator. All discrete from inside. All
analogue from outside.}

Die BAWIM ist diese Beschreibung in Reinform -- ohne Photonik, ohne
Kryogenik. Die Zustände sind Phasen 0 und $\pi$; die Anregung ist ein
kontinuierlicher RF-Träger; die Kopplung ist ein analoger Strom mit
15-bit-Auflösung. Der Spin ist kein Objekt, sondern eine Phasenlage. \SM

% ============================================================
\section{Einordnung 2: Barkhausen als Wicklungsquantisierung (Dok.~314)}
% ============================================================

Die Barkhausen-Bedingung \enquote{integer multiple of $2\pi$} ist eine
Windungszahl: Nur Zustände, bei denen die akkumulierte Phase um die
geschlossene Schleife $2\pi n$, $n\in\mathbb{Z}$, beträgt, zirkulieren
stabil. \BM\ (Dok.~314, Ch.~D2): \enquote{winding numbers are topological};
die Modenstruktur auf $T^4$ ist durch ganzzahlige Wicklung um jeden Zyklus
festgelegt.

Der Ringoszillator ist ein $S^1$ -- ein einzelner Zyklus. Die Barkhausen-%
Bedingung ist die Wicklungsquantisierung auf diesem $S^1$. Die BAWIM
realisiert damit technisch, was FFGFT für $T^4$ setzt: Zustände existieren
nur auf ganzzahligen Wicklungen. \SM

% ============================================================
\section{Einordnung 3: Parametrische $\mathbb{Z}_2$-Selektion (Dok.~339)}
% ============================================================

Die PSA wird bei $2\omega$ gepumpt und lässt nur die Phasen $0$ und $\pi$
durch. Das ist eine Periodenverdopplung: Der Zustandsraum $S^1$ der Phase
wird auf zwei Punkte $\{0,\pi\}\cong\{+1,-1\}$ reduziert.

\BM\ (Dok.~339, §3.1): Die Fixpunkte $\{+1,-1\}$ des Frobenius auf
$GF(27)^*$ sind der massive $\mathbb{Z}_2$-Sektor. \BM\ (Dok.~336): Die
Normklassen $\{+1,-1\}$ trennen leicht/schwer.

Die parametrische Selektion der BAWIM ist eine technische Realisierung
derselben Zweiteilung: Aus einem kontinuierlichen Phasenraum werden
genau zwei Fixpunkte herausgefiltert. Dass es \emph{zwei} sind und nicht
drei, liegt am Pumpen bei $2\omega$; ein Pumpen bei $3\omega$ ergäbe eine
$\mathbb{Z}_3$-Selektion $\{0,2\pi/3,4\pi/3\}$ -- die Potts-Maschine, die
Whitehead et al.\ (Ref.~12 in Vadde et al.) mit CMOS gebaut haben. \SM

% ============================================================
\section{Einordnung 4: Kopplungsregime mit Chaosschwelle (Dok.~328)}
% ============================================================

Dok.~328 behandelt Kopplungsregime gekoppelter Oszillatoren: Adler-Modell,
Synchronisation, Phase-Locking, und die Grenze, jenseits derer die Kopplung
den Oszillator zerstört. Die BAWIM setzt genau diese Grenze operativ:
Einspeisung auf 5--30\,\% des zirkulierenden Signals, \enquote{avoiding
chaotization}.

Die Existenz einer \emph{optimalen} Kopplungsstärke pro Probleminstanz
(Suppl.\ Note~5) ist der Adler-Sweet-Spot: zu schwach -- keine
Synchronisation, das System bleibt in Zufallszuständen; zu stark --
Chaotisierung, die Phasenbinarisierung bricht zusammen. Dazwischen liegt
das Fenster, in dem die Oszillatorbank ins Energieminimum fällt. \SM

% ============================================================
\section{Einordnung 5: Der Auflösungsboden (Dok.~343 §G$'$)}
% ============================================================

Dies ist die quantitativ deutlichste Einordnung. Dok.~343 §G$'$ formuliert:
Ein resonanzbasierter Rechner ist gedeckelt bei $N_{\max}\approx1/\varepsilon$,
wobei $\varepsilon$ die Kopplungstoleranz ist; er findet Minima bis zur
Toleranz, aber nicht die exakte Null. \BM

Die BAWIM zeigt dieses Verhalten in drei Ausprägungen:

\paragraph{MAX-CUT.} 95\,\% der Ising-Energie, 99,99\,\% des Schnittwerts.
Das Optimum wird bis auf Bruchteile eines Prozents erreicht -- aber nicht
exakt.

\paragraph{Number Partitioning.} Die Differenz $E=|\sum_{A}a_i-\sum_{B}a_i|$
wird auf $<0{,}1\,\%$ der Gesamtsumme gebracht, mit 100\,\% Erfolg. Die
\emph{exakte} Partition $E=0$ dagegen wird bei 32 Zahlen noch zu $\sim$90\,\%
gefunden, bei 2048 Zahlen praktisch nie (Abb.~5b). Die Kopplungsauflösung
ist 15~bit; die Zahlen liegen bei $\le178$; die Gesamtsumme bei 2048 Zahlen
liegt bei $\sim1{,}8\times10^5$. Eine Toleranz von 0,1\,\% entspricht
$\sim180$ -- das ist die Größenordnung einer einzelnen Zahl. Die Maschine
löst das Problem bis auf ein Element, nicht bis auf null.

\paragraph{Sudoku.} Bei 99,42\,\% der Grundzustandsenergie ist die Lösung
noch ungültig. Sudoku ist ein Constraint-Satisfaction-Problem: Es gibt kein
\enquote{fast richtig}. Vadde et al.\ nennen es selbst: \enquote{a
near-ground-state Ising energy may still correspond to an invalid solution.}
Dass die BAWIM das Sudoku dennoch löst, liegt an der Kleinheit des Problems
(729 Spins) und daran, dass die Grundzustandsenergie hier ein isoliertes
Minimum ist; HbSB scheitert am selben Problem.

\paragraph{Die Einordnung.} Alle drei Befunde sind konsistent mit
Dok.~343: Der Resonanzrechner ist ein Analogrechner mit Auflösungsboden.
Er ist \emph{nicht} schlechter als ein digitaler Algorithmus -- er ist bei
Näherungen besser (NPP 99,9\,\%: BAWIM schlägt HbSB durchgehend). Er ist
\emph{anders}: Er findet die Umgebung des Minimums, nicht den Punkt. Für
Optimierung ist das die richtige Aufgabe; für exakte Arithmetik nicht. \SM

% ============================================================
\section{Was nicht übereinstimmt}
% ============================================================

\paragraph{Klassisch, nicht quantenmechanisch.} Die BAWIM ist ein
klassischer RF-Ringoszillator bei Raumtemperatur. Sie beansprucht keinen
Quantenvorteil und erhebt keinen. Die FFGFT-Einordnungen betreffen die
Resonatorstruktur, nicht Quantenkohärenz.

\paragraph{Ein Zyklus, nicht vier.} Der Ringoszillator ist $S^1$; $T^4$
hat vier unabhängige Zyklen. Die Wicklungsparallele (Einordnung~2) gilt für
einen Zyklus; die $T^4$-spezifischen Aussagen (Fixpunkte, $\mathbb{Z}_3$-%
Trialität, Massenzuordnung) haben in der BAWIM kein Gegenstück.

\paragraph{Benchmark-Vorbehalt.} Vadde et al.\ betonen zweimal, dass
weder für HbSB noch für BAWIM eine Parameteroptimierung durchgeführt
wurde. Die quantitativen Vergleiche (341~ms gegen 102~ms; NPP-Erfolgsraten)
sind daher Momentaufnahmen. Die \emph{qualitative} Struktur -- Näherung
zuverlässig, Exaktheit skalenabhängig -- ist davon nicht betroffen.

\paragraph{Der Auflösungsboden ist hier technisch, nicht fundamental.}
Die 15-bit-Kopplungsauflösung ist eine Wahl der Elektronik (AD9708 DAC),
kein Naturgesetz. Dok.~343 §G$'$ spricht von $\varepsilon=\xi$ als
natürlichem Boden; die BAWIM liegt mit $2^{-15}\approx3\times10^{-5}$
zufällig in derselben Größenordnung wie $\xi=1{,}33\times10^{-4}$, ohne
dass hier ein Zusammenhang behauptet wird.

% ============================================================
\section{Zusammenfassung}
% ============================================================

\begin{center}
\begin{tabular}{p{4.4cm}p{4.4cm}p{2.6cm}l}
\toprule
BAWIM (Vadde et al.) & FFGFT & Dok. & Status \\
\midrule
Spin = Phase 0/$\pi$ eines Pulses & Ising-Maschine = Resonator & 173 & \SM \\
Barkhausen: Phase $=2\pi n$ & Wicklungsquantisierung auf $S^1$ & 314 & \SM \\
PSA bei $2\omega$: $\{0,\pi\}$ & $\mathbb{Z}_2$-Fixpunkte $\{+1,-1\}$ & 339, 336 & \SM \\
Kopplung 5--30\,\%, Chaosgrenze & Adler-Regime, Kopplungsschwelle & 328 & \SM \\
99,9\,\% ja, exakt nein & Auflösungsboden $N_{\max}\approx1/\varepsilon$ & 343 §G$'$ & \SM \\
\bottomrule
\end{tabular}
\end{center}

Die BAWIM ist der bislang klarste veröffentlichte Fall eines
Resonanzrechners im Sinne von Dok.~173/343: thermisch stabil, ohne
Nachselektion, ohne Kryogenik, mit reproduzierbaren Näherungslösungen und
einem messbaren Auflösungsboden bei exakten Lösungen. Sie bestätigt die
FFGFT-Einordnung nicht als Theorie, sondern als Klassifikation: Was
resoniert, findet Minima -- und findet sie bis zur Toleranz.

\paragraph{Quelldokument.} Vadde,~V.; Ovcharov,~R.; González,~V.~H.;
Khymyn,~R.; Litvinenko,~A.; Åkerman,~J.: \textit{A 2048-spin bulk acoustic
wave Ising machine for number partitioning and Sudoku.}
Commun.\ Phys.~\textbf{9}, 290 (2026). DOI: 10.1038/s42005-026-02846-7.
Open Access, CC-BY~4.0. Eingegangen 2.~März 2026, angenommen
20.~August 2026, veröffentlicht 8.~September 2026.
